% Rayado de 6 mm con columnas punteadas (A4) — Eric Lucero González — https://ericlucero.dev
% Licencia CC BY-NC-SA 4.0. Compilar con pdfLaTeX (dos pasadas: latexmk lo hace solo).
\documentclass[a4paper]{article}
\usepackage[T1]{fontenc}
\usepackage[spanish]{babel}
\usepackage{lmodern}
\usepackage{tikz}
\pagestyle{empty}

% ── Configuración (mm) ─────────────────────────────────────
\def\anchoPapel{210}
\def\altoPapel{297}
\def\margenX{12}     % margen mínimo a izquierda y derecha
\def\margenY{15}     % margen mínimo arriba y abajo
\def\paso{6}        % distancia entre líneas y entre columnas

\begin{document}
\begin{tikzpicture}[
    remember picture, overlay, shift={(current page.south west)},
    x=1mm, y=1mm,
    hline/.style={gray!55, line width=0.35pt},
    vline/.style={gray!55, line width=0.45pt, line cap=round, dash pattern=on 0pt off 1.5mm}
  ]
  % Celdas enteras que caben; la rejilla queda centrada en la hoja
  \pgfmathtruncatemacro{\nx}{(\anchoPapel - 2*\margenX) / \paso}
  \pgfmathtruncatemacro{\ny}{(\altoPapel  - 2*\margenY) / \paso}
  \pgfmathsetmacro{\xIni}{(\anchoPapel - \nx*\paso) / 2}
  \pgfmathsetmacro{\yIni}{(\altoPapel  - \ny*\paso) / 2}
  \pgfmathsetmacro{\xFin}{\xIni + \nx*\paso}
  \pgfmathsetmacro{\yFin}{\yIni + \ny*\paso}

  % Líneas horizontales continuas
  \foreach \i in {0,...,\ny} {
    \draw[hline] (\xIni, \yIni + \i*\paso) -- (\xFin, \yIni + \i*\paso);
  }
  % Columnas punteadas (un punto cada 1,5 mm)
  \foreach \j in {0,...,\nx} {
    \draw[vline] (\xIni + \j*\paso, \yIni) -- (\xIni + \j*\paso, \yFin);
  }
\end{tikzpicture}
\end{document}
